% !TeX program = xelatex
% Compile from this directory: latexmk -xelatex template.tex
% All paths are relative to this template directory.
\documentclass[12pt,a4paper]{article}
\usepackage[margin=2.5cm,headheight=26pt,headsep=12pt]{geometry}
\usepackage{fontspec}
% Times New Roman is the required submission font. The fallback supports TeX-only systems.
\IfFontExistsTF{Times New Roman}{\setmainfont{Times New Roman}}{%
  \setmainfont{texgyretermes-regular.otf}[BoldFont=texgyretermes-bold.otf,ItalicFont=texgyretermes-italic.otf,BoldItalicFont=texgyretermes-bolditalic.otf]
  \PackageWarningNoLine{report-template}{Times New Roman unavailable. Install it for submission}}
\usepackage{amsmath,siunitx,graphicx,booktabs,tabularx}
\usepackage{caption,subcaption,enumitem,titlesec,fancyhdr}
\usepackage[hidelinks]{hyperref}
\graphicspath{{fig/}}
\hypersetup{pdftitle={Laboratory Report Template},pdfauthor={},pdfsubject={Reusable laboratory report template}}
\setlength{\parindent}{2em}
\setlength{\parskip}{0pt}
\AtBeginDocument{\fontsize{12pt}{20pt}\selectfont}
\setlist[enumerate]{label=(\arabic*),leftmargin=2.2em,itemsep=0pt,topsep=4pt}
\titleformat{\section}{\bfseries\fontsize{14pt}{20pt}\selectfont}{\thesection.}{0.6em}{}
\titleformat{\subsection}{\bfseries\fontsize{12pt}{20pt}\selectfont}{\thesubsection}{0.6em}{}
\titlespacing*{\section}{0pt}{14pt}{6pt}
\titlespacing*{\subsection}{0pt}{8pt}{4pt}
\captionsetup{font={small},labelfont=normalfont,labelsep=period}
\captionsetup[figure]{name=Fig.,position=bottom}
\captionsetup[table]{position=top}
\pagestyle{fancy}
\fancyhf{}
% Header wording follows the supplied departmental PDF.
\fancyhead[C]{\fontsize{10.5pt}{12pt}\selectfont\bfseries
 Southern University of Science and Technology of China\\
 Department of Material Science and Technology}
\fancyfoot[C]{\thepage}
\renewcommand{\headrulewidth}{0pt}

% ---------- Fill in these fields ----------
\newcommand{\CourseName}{Course Name}
\newcommand{\ExperimentNumber}{X}
\newcommand{\ExperimentTitle}{Experiment Title}
\newcommand{\StudentName}{Student Name}
\newcommand{\StudentID}{Student ID}
\newcommand{\ExperimentDate}{YYYY-MM-DD}
\newcommand{\SubmissionDate}{YYYY-MM-DD}
\newcommand{\TeamMembers}{Team members, if any}
\newcommand{\ClassNumber}{Class number}
\newcommand{\ListNumber}{Name list number}
\begin{document}

% ==============================================================================
%                           封面部分开始 / COVER START
% ==============================================================================
\hypersetup{pageanchor=false}
\begin{titlepage}
\thispagestyle{fancy}
\fancyfoot[C]{}
% Cover layout and field order follow Lab report template.docx.
\fancyhead[C]{\fontsize{10.5pt}{12pt}\selectfont\bfseries
 South University of Science and Technology of China\\
 Department of Materials science and Engineering}
\vspace*{1.1cm}
\begin{center}
{\fontsize{18pt}{24pt}\selectfont\bfseries \CourseName\par}
\vspace{0.8cm}
{\fontsize{18pt}{24pt}\selectfont\bfseries
 Lab \ExperimentNumber: \ExperimentTitle\par}
\end{center}
\vspace{3cm}
\noindent\hfill
\begin{minipage}{0.58\textwidth}
\raggedright
\fontsize{14pt}{22pt}\selectfont\bfseries
Class No.: \ClassNumber\par
Name list number: \ListNumber\par
\vspace{1.3cm}
Name: \StudentName\par
Student ID: \StudentID\par
Team member: \TeamMembers\par
Date of experiment: \ExperimentDate\par
Date of submission: \SubmissionDate\par
\end{minipage}
\vfill
\end{titlepage}
\hypersetup{pageanchor=true}
\setcounter{page}{1}

% ==============================================================================
%                           封面部分结束 / COVER END
% ==============================================================================

\section{Objectives}
\begin{enumerate}
\item State the theory or physical principle to be examined.
\item Identify the preparation, measurement, or analysis skills to be acquired.
\item State the sample or performance target specified in the experiment handout.
\end{enumerate}

\section{Principles}
\subsection{Background and governing relation}
Explain the essential scientific background in your own words. Define the
quantities used in the analysis and cite the sources of equations or borrowed
ideas. Keep the complete principles section within one page, including figures.
Equation~\eqref{eq:example} demonstrates an editable, numbered equation.
\begin{equation}
 y = ax+b.
 \label{eq:example}
\end{equation}
Here $x$ is the independent variable, $y$ is the measured response, $a$ is the
slope, and $b$ is the intercept. Specify the units of each quantity for the actual
experiment and replace this illustrative relation with the governing equation.

\subsection{Experimental principle or setup}
Replace placeholder A in Fig.~\ref{fig:principle} with a labelled schematic of
the experimental principle or setup. Explain the components and signal path.
\begin{figure}[htbp]
 \centering
 \includegraphics[width=0.75\linewidth]{figA.jpg}
 \caption{Placeholder for a principle or setup schematic. Replace image A
 with the experimental schematic and describe its components.}
 \label{fig:principle}
\end{figure}
\clearpage

\section{Experimental materials and setups}
\subsection{Materials}
List the chemicals, samples, and consumables with their purity, concentration,
and dimensions. Use correct chemical subscripts and units, for example H$_2$O
and \SI{20}{\milli\liter}. Replace these examples with your materials.

\subsection{Apparatus and tools}
List the instruments, model numbers, ranges, and tools. State the calibration
or resolution when it affects the results.

\section{Experimental procedures}
\subsection{Sample preparation}
\begin{enumerate}
\item Describe the initial preparation and sample labels.
\item Record the quantities, temperatures, durations, and settings actually used
in class, including changes to handout parameters.
\item Describe subsequent treatment, storage, and observations in sequence.
\end{enumerate}

\subsection{Measurement and data processing}
\begin{enumerate}
\item Describe instrument setup, calibration, and the measurement conditions.
\item State the measurement sequence, repeat count, and how raw data were saved.
\item Explain data correction, fitting, and uncertainty estimation. Summarize the
sequence in Fig.~\ref{fig:workflow}.
\end{enumerate}
\begin{figure}[htbp]
 \centering
 \includegraphics[width=0.60\linewidth]{figB.jpg}
 \caption{Placeholder for an experimental workflow. Replace image B with the
 actual operations and processing conditions.}
 \label{fig:workflow}
\end{figure}
\clearpage

\section{Results and discussion}
\subsection{Presentation of results}
Organize the measured results into a three-line table such as
Table~\ref{tab:results}. Place the table number and a concise title above the
table. Give the symbol and unit of every physical quantity. Transcribe raw
records into a typeset table and retain the original records for the submission
appendix. Replace all bracketed entries with measured values.
\begin{table}[htbp]
 \centering
 \caption{Example structure for measured results. Bracketed entries are placeholders.}
 \label{tab:results}
 \begin{tabularx}{\linewidth}{@{}Xccc@{}}
 \toprule
 Sample & Time $t$ / \si{\second} & Response $y$ / \si{\volt} &
 Uncertainty $u_y$ / \si{\volt}\\
 \midrule
 Sample 1 & [value] & [value] & [value]\\
 Sample 2 & [value] & [value] & [value]\\
 Sample 3 & [value] & [value] & [value]\\
 \bottomrule
 \end{tabularx}
\end{table}

Figure~\ref{fig:results} demonstrates a two-panel layout. Use readable axis
labels with symbols and units. Explain error bars and the number of repeats in
the caption. For photographs or micrographs, include a calibrated scale bar and
magnification. Replace the letter placeholders C and D with experimental plots
or images and update the captions accordingly.
\begin{figure}[htbp]
 \centering
 \begin{subfigure}[t]{0.49\linewidth}
  \centering
  \includegraphics[width=\linewidth]{figC.jpg}
  \caption{Placeholder C for a result plot.}
  \label{fig:plot}
 \end{subfigure}\hfill
 \begin{subfigure}[t]{0.49\linewidth}
  \centering
  \includegraphics[width=\linewidth]{figD.jpg}
  \caption{Placeholder D for an experimental image.}
  \label{fig:image}
 \end{subfigure}
 \caption{Placeholders for a two-panel result figure. Replace C and D with
 experimental results. Describe the error bars, sample size, annotations,
 and imaging conditions relevant to the completed figure.}
 \label{fig:results}
\end{figure}

\subsection{Interpretation and comparison}
Describe the main trends and compare the measured values with the theoretical
prediction or expected behaviour. Use quantitative comparisons where possible.
Explain discrepancies and unexpected observations using the experimental
conditions and relevant scientific principles. Address the discussion topics
specified in the experiment handout.

\subsection{Error analysis}
Identify systematic and random sources of error, explain their effect on the
results, and estimate their magnitude where possible. Describe practical steps
to reduce the dominant errors. State the limitations that affect the conclusion.

\section{Conclusion}
Replace this paragraph with a concise summary of the experiment and its main
findings. State the method, principal quantitative results, comparison with the
objectives or theory, and the supported conclusion. Use one or two paragraphs.
The departmental guidance gives an approximate length of 300 characters, so
keep the English conclusion similarly concise and follow any course-specific limit.

\section{Questions}
\subsection*{(1) Replace with the full question from the handout}
\noindent\textbf{Answer}\quad
Explain the reasoning and show calculations where needed. Follow the question
order in the handout. Verify and summarize any AI-assisted material before use.

\section{References}
Use a consistent citation style and cite sources in the text, for example
Ref.~\cite{replace-source}. Replace the illustrative entry below with a source
actually consulted. Remove this section and its citations if no references are used.
\begingroup
\renewcommand{\section}[2]{}% The numbered References heading is already present.
\begin{thebibliography}{9}
\bibitem{replace-source}
[Author(s)]. [Title]. [Journal or publisher], [year], [volume], [pages or article number].
[DOI or URL, if applicable].
\end{thebibliography}
\endgroup

% Submission attachments. Uncomment after placing your own PDFs in this folder.
% The departmental PDF requires the report, pre-lab report, and raw records in one PDF.
% Add \usepackage{pdfpages} to the preamble before enabling the following lines.
% \clearpage
% \includepdf[pages=-]{prelab.pdf}
% \includepdf[pages=-]{raw-records.pdf}
\end{document}
